% Separate post-CP24 candidate; depends on the fully written B67 candidate.
\subsection{The initial tame ramification sector in height two}

Fix \(p\ge5\) and
\[
 e=t(p-1),\qquad 2\le t\le p-1,\qquad
 q=p^2,\quad N=q-1,\quad s=e(p+1)=tN .
 \tag{B-RD68.1}
\]
In particular \(p\nmid e\), \(e<p^2\), and \(e\) is even.
After completed maximal unramified base change, the field \(K\)
is complete of mixed characteristic \((0,p)\), with algebraically
closed residue field, \(e_K=e\) and \(\mu_p\subset K\).
For whole potentially supersingular height-two groups we take
the common good tame extension \(L/K\) of degree \(N\), and normalize
\(v_L(\pi_L)=1\). Thus \(v_L(p)=eN=(p-1)s\).

\begin{lemma}[Unequal-Newton rules throughout the sector]
\label{lem:B-RD68-rules}
\status{Candidate; all new-index valuation comparisons explicit.}
Suppose the common finite determinant character is nonsquare.
Make a common finite height-one twist so that every actual
special tangent exponent satisfies \(c\equiv1\pmod{p-1}\).
If the residual module is reducible, its Newton polygon has two
slopes, with unique high line \(H\) and integer \(B\) satisfying
\[
 \begin{gathered}
 0<B<pe,\qquad B\equiv c\pmod{p+1},\qquad B\ \text{odd},\\
 \alpha=s-B,\quad\beta=B/p,\quad\alpha>\beta,\\
 \text{high-line character } \zeta^{c-B}.
 \end{gathered}
 \tag{B-RD68.2}
\]
Every order-\(p^2\) point whose \(p\)-multiple is outside \(H\)
has \(v_L\)-valuation \(B/p^3\).
For two such groups with isomorphic actual \(p^2\)-torsion:

\begin{enumerate}
\item If their residual wild image is nontrivial, the high lines
agree. Two parameters prime to \(p\) are equal and give equal
special characters. Exactly one parameter cannot be divisible
by \(p\).
\item If the residual image is tame, write \(B_i=pb_i\).
When both \(b_i\) are prime to \(p\), they are equal, their
high lines agree, and the special characters agree.
\end{enumerate}
\end{lemma}
\begin{proof}
The universal height-two coefficient bounds established in
Lemma~\ref{lem:B-RD49-point} give the only possible negative
Newton vertices \((1,(p-1)s),(p,v_L(u)),(q,0)\).
If all nonzero first-level points have one valuation, their
leading-coefficient map is injective. Wild inertia acts trivially
on that map: its uniformizer multiplier and residue action are
both trivial, and the good descent is tame. Thus the residual
module is tame, with a scalar leading-coordinate action.
If it is reducible over \(\mathbf F_p\), that scalar is in
\(\mathbf F_p^\times\), and its determinant character is square.
Because \(\mu_p\subset K\), reduction of the cyclotomic determinant
is trivial, so this contradicts the nonsquare finite determinant.
The reducible case must therefore have two slopes.

The high line is stable under inertia, and its character has
values in \(\mathbf F_p^\times\). It is rational over \(L\).
Its valuation is integral, so writing
\(v_L(u)=(p-1)B\) gives the integer \(B\) and the slopes
in (B-RD68.2). Their strict order gives \(B<pe\).
Choose an exact tame coordinate and
\(\sigma\pi_L=\zeta\pi_L\), as in the integral averaging proof
of Lemma~\ref{lem:B-RD67-normalization}. On a nonzero high
point the leading character is \(\zeta^{c+\alpha}\).
Since it is \(\mathbf F_p\)-valued,
\(c+s-B\equiv0\pmod{p+1}\), giving the asserted congruence.
Here \(s=tN\), so its actual value is \(\zeta^{c-B}\).
As \(c\) is odd and \(p+1\) is even, \(B\) is odd.
The inverse Newton polygon at target valuation \(\beta\)
is one segment of degree \(q\), giving \(B/p^3\) for every
preimage.

Let \(R/L\) be the common residual division field. It is trivial
or cyclic of degree \(p\). In the latter case the high line
is the unique fixed line of the wild unipotent image, so it is
common. A low point has \(v_R\)-valuation \(B\), while its
nonzero wild displacement has valuation \(p\alpha\).
The fixed-coefficient lower-break calculation gives
\[
 r=ps-(p+1)B\quad(p\nmid B),\qquad
 r<ps-(p+1)B\quad(p\mid B).
 \tag{B-RD68.3}
\]
Indeed a coordinate of valuation \(B\) changes in valuation
\(B+r\) if \(p\nmid B\), and strictly later if \(p\mid B\);
every higher term changes later still. This is precisely the
proof of Lemma~\ref{lem:B-RD60-break}, with the present valuations.
Two prime-to-\(p\) parameters give the same \(r\), hence the
same \(B\); their common high-line character then gives the
same \(c\) modulo \(N\).

If \(p\nmid B_1\) and \(p\mid B_2\), (B-RD68.3) gives \(B_1>B_2\).
Choose corresponding primitive points \(P_i\) with \(pP_i\)
outside the common high line. Their \(v_R\)-valuations are
\(B_i/p^2\). Their point fields over \(R\) are equal, because
the given Galois-module isomorphism identifies their stabilizers.
The denominator on the first side forces their common point
field \(E\) to have degree \(p^2\): it is a torsor under
the \(L\)-rational first-level module, so degree is at most \(p^2\).
Thus every first-level translation occurs.
In \(E\), the two point coordinates have valuations \(B_i\).
A translation outside the common high line has displacement
valuation \(p^2B_i\). The first coordinate gives actual break
\(NB_1\); the universal lower bound on the second gives
\[
 p^2B_2\ge B_2+NB_1,
 \tag{B-RD68.4}
\]
contradicting \(B_1>B_2\). This proves the first rule.

In the tame case the residual module is rational over \(L\),
so \(B_i=pb_i\). Choose \(pP_i\) outside both high lines.
If both \(b_i\) are prime to \(p\), the common point field has
degree \(p^2\), and the point coordinates there have valuations
\(b_i\). A translation outside both high lines gives the same
break \(Nb_1=Nb_2\), hence \(b_1=b_2=b\).
If the high lines differed, a translation high only in one
model would have unequal displacement valuations, although
both coordinates have the same prime-to-\(p\) valuation \(b\)
and must change in valuation \(b+r\). This is impossible.
The high-line characters now identify \(c\) modulo \(N\).
\end{proof}

\begin{lemma}[The actual cyclic quotient transformations]
\label{lem:B-RD68-quotients}
\status{Candidate; integral maps, root multiplicities and torsion levels specified.}
The actual quotient by the high line has Newton parameter
\[
 B_H=\begin{cases}
       pB,&B<e,\\
       s-B,&B>e.
      \end{cases}
 \tag{B-RD68.5}
\]
There is no equality case, because \(B\) is odd and \(e\) even.
If the residual module is rational over \(L\) and \(J\) is an
inertia-stable low line, its actual quotient has parameter
\[
 B_J=B/p .
 \tag{B-RD68.6}
\]
Each quotient is good over \(\mathcal O_L\), descends generically
to \(K\), and has special tangent exponent \(pc\) modulo \(N\).
Its finite special character is unchanged. A \(p^3\)-torsion
isomorphism preserving the chosen lines gives a \(p^2\)-torsion
isomorphism of these quotients.
\end{lemma}
\begin{proof}
The finite flat closure of a chosen rank-\(p\) line has an
actual good formal quotient. Its coordinate map is a unit times
\(\prod_{Q\in J}(T-T(Q))\). Valuations are therefore computed
by this product, without assuming that a rational isogeny
preserves a chosen parameter thickness.

For \(J=H\), every point below valuation \(\alpha\) has its
valuation multiplied by \(p\). The image of \(M[p]/H\) has
valuation \(B\). For a nonzero high target, if \(B>e\) the
inverse polygon is one segment with \(q\) roots of valuation
\(\alpha/q\); the resulting \(q-p\) nonzero quotient points
have valuation \(\alpha/p\). The other \(p-1\) points have
valuation \(B>\alpha/p\), giving \(B_H=s-B\).
If \(B<e\), each nonzero high target has \(p\) inverse roots
of valuation \((s-pB)/p\) and \(q-p\) roots of valuation \(B/p\).
All lie below \(\alpha\). After quotienting the root counts
by \(p\), and including \(M[p]/H\), there are \(p-1\)
points of valuation \(s-pB\) and \(q-p\) of valuation \(B\).
Thus \(B_H=pB\).

For a low line \(J\), every nonzero class in \(M[p]/J\)
has one high representative. Its quotient valuation is
\(\alpha+(p-1)\beta=s-B/p\). All preimages of nonzero low
targets have valuation \(\beta/q<\beta\), and their quotient
valuations are \(\beta/p=B/p^2\). Their count is \(q-p\).
This gives (B-RD68.6); \(B/p\) is integral because the
first-level module is \(L\)-rational.

On the special fibre the unique subgroup of rank \(p\) is
the Frobenius kernel. Thus the special quotient is the
Frobenius twist and has inherited tangent exponent \(pc\).
The unordered pair \(\{c,pc\}\) is unchanged.
The chosen line is stable under the actual generic descent,
so the generic quotient is defined over \(K\).
Finally its \(p^2\)-torsion is exactly
\[
 \{P\in M[p^3]:p^2P\in J\}/J .
 \tag{B-RD68.7}
\]
This identity of actual finite generic groups proves the
claimed passage between the two torsion levels.
\end{proof}

\begin{proposition}[Nonsquare formal comparisons in the initial sector]
\label{prop:B-RD68-nonsquare}
\status{Candidate; complete residual and quotient case split written.}
Two whole potentially supersingular height-two groups in this
sector with common nonsquare finite determinant and isomorphic
actual \(p^3\)-torsion have equal finite special inertia characters.
\end{proposition}
\begin{proof}
Normalize the common finite determinant by a common finite
height-one twist. An irreducible residual comparison is already
settled at \(p^2\) by Theorem~\ref{thm:B-RD67}.
The other comparisons have the unequal parameters just proved.

In the residual wild case, their high line is common.
The rules handle any case with a prime-to-\(p\) parameter.
Suppose both \(B_i\) are divisible by \(p\).
If both are below \(e\), write \(B_i=pb_i\).
Their common high-line character and \(c_i\equiv1\pmod{p-1}\)
give \(b_1\equiv b_2\pmod{p-1}\). But
\[
 0<b_i<e/p\le (p-1)^2/p<p-1,
 \tag{B-RD68.8}
\]
so \(b_1=b_2\), and the high-line character identifies \(c_1=c_2\).
If both \(B_i>e\), quotient by the common high line.
The new parameters \(s-B_i\) are prime to \(p\), since \(p\nmid s\).
Their actual \(p^2\)-torsion is isomorphic, so the rules identify
the quotient special characters, and therefore the original ones.
If one parameter is below \(e\) and one above, the two quotient
parameters \(pB_i\) and \(s-B_j\) have opposite \(p\)-divisibility.
They cannot have isomorphic \(p^2\)-torsion by the same rules.

Now suppose residual inertia is tame, and write \(B_i=pb_i\),
where \(0<b_i<e\). Two \(b_i\) prime to \(p\) are already
handled. Suppose \(p\nmid b_1=b\) and \(b_2=pk\).
Then \(B_2=p^2k>e\), since \(e<(p^2)\).
Choose corresponding primitive points whose \(p\)-multiples
are outside both high lines. The common point field has degree
\(p^2\) over \(L\), forced by the first coordinate valuation
\(b/p^2\); its normalized coordinate valuations are \(b\) and \(pk\).
If the high lines differ, choose a translation high in model 1
and low in model 2. The first coordinate, prime to \(p\), fixes
its actual lower break as \(p^2\alpha_1-b\).
The universal lower bound in model 2 would require
\[
 p^2(\alpha_1-pk)\le b-pk .
 \tag{B-RD68.9}
\]
Yet \(\alpha_1=s-pb\ge e+p\), while \(pk<e\), \(b<e<p^2\).
The left side is greater than \(p^3\), and the right side less
than \(p^2\), a contradiction. Thus the high lines agree.
Quotient by this common line. The second parameter \(s-B_2\)
is prime to \(p\). If \(B_1<e\), the first parameter \(pB_1\)
is divisible, excluded by the rules. If \(B_1>e\), both new
parameters are prime to \(p\), so their equality forces
\(B_1=B_2\), contradicting \(p\nmid b\).

Finally let \(b_i=pk_i\) for both groups. Both
\(B_i=p^2k_i>e\). If the high lines agree, their high quotients
have prime-to-\(p\) parameters and equal characters by the rules.
If the high lines differ, quotient the first group by its
high line and the second by the corresponding stable low line.
The parameters are \(s-B_1\), prime to \(p\), and
\(B_2/p=pk_2\), divisible by \(p\). Their inherited actual
\(p^2\)-torsion isomorphism is impossible. These are all
residual and divisibility cases.
\end{proof}

\begin{lemma}[Square determinants in the same sector]
\label{lem:B-RD68-square}
\status{Candidate; the coefficient bound is imposed on the actual target.}
The same conclusion holds when the common finite determinant
character is square.
\end{lemma}
\begin{proof}
Twist both groups by its common finite square root.
Their special inertia eigenvalues have norm one, so their
simultaneous good field has degree \(m\mid p+1\).
Put \(a=e m\le e(p+1)\).
If the special characters differ, their residual tame spectra
are disjoint. The integral Tate pairings and self-orthogonal
graph argument of Proposition~\ref{prop:B-RD64-formal} give
the special graph closure
\(\ker F_M^3\times\ker F_N^3\).
In detail the graph is isotropic generically, hence integrally
by flatness; equality with its orthogonal follows from equal
finite flat ranks. Integral tame polynomial projectors split
its special fibre into the two formal factors, and uniqueness
of rank-\(p^a\) subgroups of a one-dimensional connected formal
group gives exponent three on each factor.
The actual quotient \(Q=(M\times N)/H\) consequently has
special fibre \(M_0^{(p^3)}\times N_0^{(p^3)}\).
The map
\[
 M\times N\longrightarrow Q\times M,\qquad
 (x,y)\longmapsto(q(x,y),x)
\]
has no nonzero finite flat subgroup in its kernel, while its
special kernel is \(0\times\ker F_N^3\).

The actual target is a product of height-two formal groups
on the special fibre. For a coefficientwise valuation-one
vector its first two \(p\)-multiplications have valuations
\[
 q_1\ge p+1,\qquad
 q_2\ge\min\{a+p+1,p^2+p+1\}\ge a/(p-1).
 \tag{B-RD68.10}
\]
Here \(a\ge e\ge2(p-1)\ge p\).
The differential identity for the \(p\)-series makes any
coefficient with a non-\(p\)-divisible exponent divisible by \(p\);
all total degrees below \(p^2\) vanish specially, so the
remaining such coefficients are divisible by a uniformizer.
These are precisely the intrinsic coefficient estimates on
the actual target used in B34 and B64.
Finally
\[
 a/(p-1)\le e(p+1)/(p-1)\le(p-1)(p+1)=N .
 \tag{B-RD68.11}
\]
The closed coefficientwise logarithm and Cartier bounded-projection
and saturation proof in Lemma~\ref{lem:B-RD64-weighted-kernel},
using the just stated unweighted bounds, forces \(F^2\) to
kill this special kernel. It cannot kill \(\ker F_N^3\).
This contradiction uses the actual quotient target, without
transferring any Hasse-coefficient thickness from \(M\) or \(N\).
\end{proof}

\begin{theorem}[The exact initial tame sector]
\label{thm:B-RD68}
\status{Candidate; complete full-monodromy reduction and sharp lower examples written.}
For every prime \(p\ge5\), in the full B.4 monodromy setting
with height two, perfect residue field and \(\mu_p\subset K\),
\[
 \boxed{\displaystyle
 \nu_{\min}\bigl(p,t(p-1),2\bigr)=2
 \quad\text{for every }2\le t\le p-1.}
 \tag{B-RD68.12}
\]
Thus actual \(p^3\)-torsion always determines the requested
character in this sector, whereas actual \(p^2\)-torsion
does not do so uniformly. In particular
\(\nu_{\min}(5,16,2)=2\).
\end{theorem}
\begin{proof}
For the lower bound, \(e=t(p-1)>p\) and \(p-1\mid e\).
The explicit actual \(p^2\)-torsion comparison in
Theorem~\ref{thm:B-RD63-arc} therefore applies at every
one of these indices.

For the upper bound, pass to the completed maximal unramified
base and compare actual \(p^3\)-torsion. The determinant of
an allowed height-two monodromy system is
\(\chi_p^d\epsilon\), \(d\in\{0,1,2\}\), with \(\epsilon\)
finite and \(\mu_{p-1}\)-valued. Since \(p\nmid e\), the
cyclotomic image modulo \(p^3\) has order \(p^2\).
Equality of determinants gives the same \(d\), because
\(|d_1-d_2|\le2<p\), and then the same \(\epsilon\).
For \(d=0\), the special character is the Brauer lift of the
semisimplified residual inertia character; \(d=2\) follows
by Cartier duality. For \(d=1\), each group is either whole
potentially supersingular or filtered by two potentially good
height-one groups. Two filtered systems have their finite
characters determined modulo \(p\), since \(\mu_p\subset K\).
A mixed comparison is impossible by
Proposition~\ref{prop:B-RD63-mixed}, because \(e<p^2<p^3\).
The two whole formal groups have the common finite determinant
just established. Propositions~\ref{prop:B-RD68-nonsquare}
and Lemma~\ref{lem:B-RD68-square} settle its two square classes.
This proves the upper bound, and hence the exact index.
\end{proof}

\begin{remark}[Boundary of the sector]
\label{rem:B-RD68-scope}
\status{Candidate scope; no assertion beyond the displayed sector.}
The proof uses \(e<p^2\) and \(e/p<p-1\) in the reducible
residual comparisons, and the actual-target estimate
\(e(p+1)/(p-1)\le p^2-1\) for square determinants.
It does not extrapolate those inequalities to \(t\ge p\)
or ramification divisible by \(p\).
The minimal index \(t=1\) has the separately proved value one
from B61. At \(p=5\), the new sector consists of \(e=8,12,16\);
the first two were already proved in B64 and B65.
The theorem concerns the full filtered monodromy condition,
not only comparisons of two whole potentially good groups.
The retained exact supplement
\path{code/reduction_descent/verify_initial_tame_sector.py}
checks all \(60\) sectors \(2\le t\le p-1\) at
\(p=5,7,11,13,17,19\). It independently computes
\(56{,}916\) Newton/high-quotient records, \(3{,}420\) low
quotients, \(113{,}832\) tangent weights, \(63{,}018\)
mixed-divisibility point-field inequalities and \(322\)
actual-target logarithm bounds, together with every stated
small-parameter and two-divisible congruence in those ranges.
The author run passed in \(2.3066381700045895\) seconds at
20:57:28 UTC. The all-prime assertion is the preceding proof
and does not follow from this finite enumeration.
\end{remark}
